
\chapter{Brauer's Theorem}
Brauer's theorem is one of the deepest results in finite group representation theory. 
 It implies that Artin L-functions are meromorphic, which we will introduce in Section 12, and provides a powerful tool for constructing irreducible representations.
\begin{defn}[Elementary]
A subgroup \(H \le G\) is called \textbf{$p$-elementary} if \(H \cong C_m \times_{grp} P\) where \(C_m\) is cyclic with order $m$ and \(P\) is a \(p\)-group (possibly trivial) for some prime $p$.
\end{defn}

\begin{remark}
Every cyclic subgroup is \(p\)-elementary: if \(C\) has order \(p^a m\) with \((p,m)=1\), then \(C = C_m \times_{grp} C_{p^a}\), where \(C_m\) is cyclic of order \(m\) and \(C_{p^a}\) is a cyclic \(p\)-group.
\end{remark}

\begin{defn}
Let \(\mathbf{R}(G)\) be the ring generated by characters of irreducible complex representations of $G$ with
\begin{itemize}
    \item addition: \((\chi_1 + \chi_2)(x) = \chi_1(x) + \chi_2(x),\quad\forall(\chi_1,\chi_2)\in \mathbf{R}(G)^2,\forall x\in G\);
    \item multiplication: \((\chi_1 \cdot \chi_2)(x) = \chi_1(x) \cdot \chi_2(x),\quad\forall(\chi_1,\chi_2)\in \mathbf{R}(G)^2,\forall x\in G\).
\end{itemize}
This corresponds to direct sum and tensor product of representations, respectively.

 The elements in $\mathbf{R}(G)$ are called \textbf{virtual characters}.

 For a fixed prime \(p\), let \(V_p \subseteq \mathbf{R}(G)\) be the ideal generated by all characters of the form \(\operatorname{Ind}_H^G \chi\) (by Proposition \ref{prop:ind eq ind res}), where \(H\) is a \(p\)-elementary subgroup of \(G\) and \(\chi\) is any character of \(H\).
\end{defn}

We first present some lemmas.

\begin{lemma}
Let \(C\) be a cyclic group of order \(n\). Define \(\theta_C: C \to \mathbb{Z}\) by
\[
\theta_C(x) =
\begin{cases}
n, & \text{if } x \text{ generates } C,\\
0, & \text{otherwise}.
\end{cases}
\]
Then \(\theta_C \in \mathbf{R}(C)\).
\end{lemma}


\begin{proof}
For each divisor \( d \mid n \), let \( C_d \) be the unique subgroup of \( C \) of order \( d \). Let
\[
\mathbf{1}_{C_d} : C_d \to \mathbb{C}^\times, \qquad x\mapsto1
\]
be the trivial representation of \( C_d \). Define
\[
\psi_d := \chi_{\operatorname{Ind}_{C_d}^{C} \mathbf{1}_{C_d}}.
\]
By Proposition \ref{Character Formula for Induced Representations}, we can get
\[
\psi_d(x) =
\begin{cases}
\dfrac{n}{d}, & \text{if } x \in C_d,\\[6pt]
0, & \text{otherwise}.
\end{cases}
\]
Define
\[
f_d (x):= \frac{d}{n} \, \psi_d (x)=
\begin{cases}
1, & \text{if } x \in C_d,\\
0, & \text{otherwise}.
\end{cases},
\]

and
\[
h_m(x) =
\begin{cases}
1, & \text{if } \operatorname{ord}(x) = m,\\
0, & \text{otherwise}.
\end{cases}
\]
Hence
\[
f_d(x) = \sum_{\substack{m \mid d \\ m \mid n}} h_m(x)=\sum_{\substack{m \mid d}} h_m(x).
\]
By Möbius inversion,
\[
h_n = \sum_{d \mid n} \mu\!\left( \frac{n}{d} \right) f_d.
\]
Then
\[
\theta_C= n \cdot h_n
= n \sum_{d \mid n} \mu\!\left( \frac{n}{d} \right) \frac{d}{n} \psi_d
= \sum_{d \mid n} d \, \mu\!\left( \frac{n}{d} \right) \psi_d.
\]
Each \( \psi_d\in \mathbf R(C) \) and the coefficients \( d \, \mu(n/d) \) are integers. Hence \( \theta_C \) is an integral linear combination of characters, so \( \theta_C \in \mathbf{R}(C) \).
\end{proof}

\begin{lemma}[Decomposition of the Constant Function \(|G|\)]\label{Decomposition of the Constant Function}
Let \(G\) be a finite group. Then for every \(s \in G\),
\[
|G| = \sum_{\text{cyclic subgroup }C \le G} \operatorname{Ind}_C^G \theta_C(s),
\]
where \(\operatorname{Ind}_C^G \theta_C\) denotes the induced character.
\end{lemma}
\begin{proof}
Fix \(s \in G\), and let \(n\) be the order of \(s\). Let \(C_s = \langle s \rangle\).

By Proposition \ref{Character Formula for Induced Representations}, for any cyclic subgroup \(C \le G\),
\[
\operatorname{Ind}_C^G \theta_C(s)
= \frac{1}{|C|} \sum_{\substack{t \in G \\ t^{-1} s t \in C}} \theta_C(t^{-1} s t).
\]
Now \(\theta_C(y) \neq 0\) only when \(y\) generates \(C\). Thus the term is nonzero only if \(t^{-1} s t\) generates \(C\), which means \(C = \langle t^{-1} s t \rangle\), a conjugate of \(C_s\).

Moreover, for a fixed cyclic subgroup \(C\) that is conjugate to \(C_s\), say \(C = t_0^{-1} C_s t_0\), every \(t \in t_0 C_s\) (a coset of \(C_s\) in \(G\)) satisfies \(t^{-1} s t \in C\) and generates \(C\). The number of such \(t\) is \(|C_s|\), and each contributes \(\theta_C(t^{-1} s t) = |C|\).

Thus the sum over all \(t \in G\) with \(t^{-1} s t \in C\) contributes \(|C_s| \cdot |C|\). Dividing by \(|C|\) gives \(|C_s| = n\).

Summing over all cyclic subgroups \(C \le G\), each conjugate of \(C_s\) contributes \(n\), and all other cyclic subgroups contribute \(0\). The number of distinct conjugates of \(C_s\) in \(G\) is \(|G|/|N_G(C_s)|\), and each contributes \(n\). But note: the sum is over all cyclic subgroups, not just one representative per conjugacy class. So the total contribution is:
\[
\sum_{\text{cyclic subgroup }C \le G} \operatorname{Ind}_C^G \theta_C(s)
= \sum_{\substack{\text{cyclic subgroup }C \le G \\ C \text{ conjugate to } C_s}} n
= \frac{|G|}{|N_G(C_s)|} \cdot n
= |G|.
\]
\end{proof}
\begin{lemma}[Class Functions with Values in \(|G|\mathbf{A}\)]
\label{Class Functions to linear combination}
Let $G$ be a finite group and let \(\mathbf{A} := \mathbb{Z}[\zeta_{|G|}]\), where \(\zeta_{|G|}\) is a primitive \(|G|\)-th root of unity. If \(f: G \to |G|\mathbf{A}\) is a class function, then \(f\) is an \(\mathbf{A}\)-linear combination of characters induced from cyclic subgroups.
\end{lemma}

\begin{proof}
Write \(f = |G| \cdot h\) with \(h: G \to \mathbf{A}\) a class function. By Lemma \ref{Decomposition of the Constant Function} and Proposition \ref{prop:ind eq ind res},
\[
f(s) = \sum_{C \le G} \operatorname{Ind}_C^G \theta_C(s) \cdot h(s)
= \sum_{C \le G} \operatorname{Ind}_C^G (\theta_C \cdot \operatorname{Res}_C^G h)(s).
\]
For each cyclic subgroup \(C\le G\), the function \(\theta_C \cdot \operatorname{Res}_C^G h\) takes values in \(|C|\mathbf{A}\). Therefore each induced term is an \(\mathbf{A}\)-linear combination of induced characters.
\end{proof}
\begin{lemma}\label{tensor product in R}
Let \(H_1, H_2\) be finite groups. If \(\chi_1 \in \mathbf{R}(H_1)\) and \(\chi_2 \in \mathbf{R}(H_2)\), then  \(\chi_1 \otimes \chi_2 \in \mathbf{R}(H_1 \times_{grp} H_2)\).
\end{lemma}

\begin{proof}
Since \(\chi_1 \in \mathbf{R}(H_1)\), it can be written as linear combination of irreducible characters (i.e., characters of irreducible representations) of \(H_1\):
\[
\chi_1 = \sum_i n_i \psi_i, \quad n_i \in \mathbb{Z}.
\]
Similarly,
\[
\chi_2 = \sum_j m_j \phi_j, \quad m_j \in \mathbb{Z}.
\]
Then
\[
(\chi_1 \otimes \chi_2)(h_1, h_2)
= \sum_{i,j} n_i m_j \, \psi_i(h_1) \phi_j(h_2).
\]
For each pair \((i,j)\), the function \((h_1, h_2) \mapsto \psi_i(h_1) \phi_j(h_2)\) is the character of the tensor product of the irreducible representations corresponding to \(\psi_i\) and \(\phi_j\), 
then by Corollary \ref{cor:Schur eq and irreducible}[2], we have\[
\frac{1}{|H_1|\,|H_2|} \sum_{h_1 \in H_1} \sum_{h_2 \in H_2} |\chi_{\psi_i}(h_1)|^2 \, |\chi_{\phi_j}(h_2)|^2
=
\left( \frac{1}{|H_1|} \sum_{h_1 \in H_1} |\chi_{\psi_i}(h_1)|^2 \right)
\left( \frac{1}{|H_2|} \sum_{h_2 \in H_2} |\chi_{\phi_j}(h_2)|^2 \right)
= 1 \cdot 1 = 1.
\]
hence $\psi_i\otimes\phi_j$ is an irreducible character of \(H_1 \times_{grp} H_2\). Thus \(\chi_1 \otimes \chi_2 \in \mathbf{R}(H_1 \times_{grp} H_2)\).
\end{proof}
\begin{defn}
Let \(\mathbf{A} = \mathbb{Z}[\zeta_{|G|}]\), and let \(p\) be a prime. We say that two elements \((a, b) \in \mathbf{A}^2\) are congruent modulo \(p\mathbf{A}\), written
\[
a \equiv b \pmod{p\mathbf{A}},
\]
if \(a - b \in p\mathbf{A}\).
\end{defn}

\begin{lemma}[Existence of a Mod \(p\) Nonzero Function in \(V_p\)]
\label{existence of nonzero function}
Let $G$ be a finite group and let \(\mathbf{A} := \mathbb{Z}[\zeta_{|G|}]\), where \(\zeta_{|G|}\) is a primitive \(|G|\)-th root of unity. There exists a class function \(\psi: G \to \mathbf{A}\) such that:
\begin{enumerate}
    \item \(\psi \in \mathbf{A}\langle V_p\rangle\);
    \item \(\psi(s) \notin{p\mathbf{A}},\forall s \in G\).
\end{enumerate}
\end{lemma}
\begin{proof}
Let \( x_1, \dots, x_r \) be representatives of the conjugacy classes consisting of elements whose order is prime to \( p \) of \( G\). For each \( i \), proceed as follows:

Let \( C_i := \langle x_i \rangle \)and let \( P_i \) be a Sylow \( p \)-subgroup of the centralizer \( C_G(x_i) \). Then
\[
H_i := C_i \times_{grp} P_i
\]
is a \( p \)-elementary subgroup of \( G \).

Define a function \( \eta_i : H_i \to \mathbf{A} \) by
\[
c p\mapsto \theta_{C_i}(c),
\]
where $(c,p)\in C_i \times P_i$. By Lemma \ref{tensor product in R}, \( \eta_i =\theta_{C_i}\otimes\mathbf 1_{P_i}\) is a virtual character of \( H_i \).

Define
\[
\psi_i := \operatorname{Ind}_{H_i}^G \eta_i.
\]
By definition, \( \psi_i \in V_p \subseteq \mathbf{R}(G) \).

We claim that, for each \( i \),
\[
\psi_i(x_j) \equiv
\begin{cases}
|C_i| \cdot \frac{|C_G(x_i)|}{|P_i|} & \text{if } i = j, \\
0 & \text{if } i \neq j,
\end{cases}
\]
in the sense of divisibility modulo \( p\mathbf{A} \). 

More precisely, let \( R_i \subset G \) be a set of representatives for the left cosets \( G/H_i \), by Proposition \ref{Character Formula for Induced Representations},
\[
\psi_i(x_j) = \sum_{\substack{r \in R_i \\ r^{-1} x_j r \in H_i}} \eta_i(r^{-1} x_j r).
\]

We now evaluate \(\psi_i(x_j)\) modulo \(p\mathbf{A}\):

\begin{itemize}
    \item If \(i \neq j\), then \(x_j\) is not conjugate to any element of \(C_i\) that generates \(C_i\), since \(x_j\) lies in a different \(p'\)-conjugacy class. Hence the sum is empty, so
    \[
    \psi_i(x_j) = 0.
    \]

    \item If \(i = j\), then \(r^{-1} x_i r \in H_i\) and \(\eta_i(r^{-1} x_i r) \neq 0\) only when \(r^{-1} x_i r = c p\) with \(c \in C_i\) generating \(C_i\). The number of such \(r \in R_i\) is exactly
    \[
    \frac{|C_G(x_i)|}{|P_i|},
    \]
    and each contributing term equals \(|C_i|\). Hence
    \[
    \psi_i(x_i) = |C_i| \cdot \frac{|C_G(x_i)|}{|P_i|}.
    \]
    Since \(|C_G(x_i)|/|P_i|\) is prime to \(p\) (as \(P_i\) is a Sylow \(p\)-subgroup of \(C_G(x_i)\)), this value is not divisible by \(p\), so
    \[
    \psi_i(x_i) \notin p\mathbf{A}.
    \]
\end{itemize}

Now define
\[
\psi := \sum_{i=1}^r \psi_i.
\]
Then \(\psi \in \mathbf{A}\langle V_p \rangle\), since each \(\psi_i \in V_p\) and the coefficients are \(1\).

It remains to prove that \(\psi(s) \notin p\mathbf{A}\) for all \(s \in G\). Let \(s \in G\), and let \(s'\) be its \(p'\)-component. Since \(s\) and \(s'\) differ by a \(p\)-element, and all character values are sums of roots of unity, the standard reduction modulo \(p\) gives
\[
\psi(s) \equiv \psi(s') \pmod{p\mathbf{A}}.
\]
The element \(s'\) is conjugate to some \(x_i\). Since \(\psi\) is a class function,
\[
\psi(s') = \psi(x_i) = \sum_{k=1}^r \psi_k(x_i).
\]
By the evaluation above, \(\psi_k(x_i) = 0\) for \(k \neq i\), and \(\psi_i(x_i) \notin p\mathbf{A}\). Hence
\[
\psi(s') \equiv \psi_i(x_i) \notin p\mathbf{A}.
\]
Thus \(\psi(s) \equiv \psi(s') \notin p\mathbf{A}\), so \(\psi(s) \notin p\mathbf{A}\) for all \(s \in G\).
\end{proof}
\begin{lemma}
Let \(V_p \subseteq \mathbf{R}(G)\) be a subgroup, and let \(\mathbf{A}\) be a commutative ring that is free as a \(\mathbb{Z}\)-module (e.g., \(\mathbf{A} = \mathbb{Z}[\zeta_n]\)). Then the natural map
\[
\mathbf{R}(G) \hookrightarrow \mathbf{A} \otimes_{\mathbb{Z}} \mathbf{R}(G)
\]
is injective, and
\[
\mathbf{A}\langle V_p \rangle \cap \mathbf{R}(G) = V_p.
\]
\end{lemma}

\begin{proof}
Since \(\mathbf{R}(G)\) is a free abelian group, say \(\mathbf{R}(G) \cong \mathbb{Z}^r\), and \(\mathbf{A}\) is free as a \(\mathbb{Z}\)-module, the tensor product \(\mathbf{A} \otimes_{\mathbb{Z}} \mathbf{R}(G) \cong \mathbf{A}^r\). The map
\[
\mathbf{R}(G) \to \mathbf{A} \otimes_{\mathbb{Z}} \mathbf{R}(G), \qquad x \mapsto 1 \otimes x,
\]
is injective because \(\mathbb{Z} \hookrightarrow \mathbf{A}\) is injective and the induced map on free modules \(\mathbb{Z}^r \to \mathbf{A}^r\) is componentwise injection.

Now let \(x \in \mathbf{A}\langle V_p \rangle \cap \mathbf{R}(G)\). Since \(x \in \mathbf{A}\langle V_p \rangle\), it can be written as
\[
x = \sum_{i=1}^n a_i v_i, \quad a_i \in \mathbf{A},\ v_i \in V_p,\forall i\in \{1,\dots,n\}.
\]
Since \(V_p \subseteq \mathbf{R}(G)\), each \(v_i\) is an integral linear combination of the \(\mathbb{Z}\)-basis of \(\mathbf{R}(G)\). Hence \(x\) is an \(\mathbf{A}\)-linear combination of basis elements of \(\mathbf{R}(G)\). But \(x \in \mathbf{R}(G)\) means its coefficients in this basis are integers. Therefore the same coefficients must be integers, so \(x\) is an integral linear combination of the \(v_i\), hence \(x \in V_p\).

Thus \(\mathbf{A}\langle V_p \rangle \cap \mathbf{R}(G) \subseteq V_p\). The reverse inclusion \(V_p \subseteq \mathbf{A}\langle V_p \rangle \cap \mathbf{R}(G)\) is immediate. Therefore equality holds.
\end{proof}
\begin{thm}[The \(p\)-Free Part Lies in \(V_p\)]
\label{m in V_p}
Let \(|G| = p^n m\) with \((p,m) = 1\). Then the constant character \(m \cdot\mathbf1_G\) lies in \(V_p\).
\end{thm}

\begin{proof}
Let \(\psi \in \mathbf{A}\langle V_p \rangle\) be the function defined in Lemma \ref{existence of nonzero function}.

Let \(N = \varphi(p^n)\) be the order of the group \((\mathbb{Z}/p^n\mathbb{Z})^\times\) , where $\varphi$ is Euler's totient function. For each \(s \in G\), since \(\psi(s) \notin p\mathbf{A}\), its residue modulo \(p\mathbf{A}\) is nonzero. Hence \(\psi(s)\) is a unit modulo \(p\mathbf{A}\), so by Euler's theorem in the ring \(\mathbf{A}/p^n\mathbf{A}\),
\[
\psi(s)^N \equiv 1 \pmod{p^n \mathbf{A}}.
\]
Thus the function
\[
s \mapsto m \big( \psi(s)^N - 1 \big)
\]
takes values in \(p^n m \mathbf{A} = |G| \mathbf{A}\).

By Lemma \ref{Class Functions to linear combination}, this function lies in \(\mathbf{A}\langle V_p \rangle\), since every cyclic subgroup is \(p\)-elementary.

Now \(\psi \in \mathbf{A}\langle V_p \rangle\), and \(\mathbf{A}\langle V_p \rangle\) is an ideal of \(\mathbf{A}\langle \mathbf{R}(G) \rangle\), so \(\psi^N \in \mathbf{A}\langle V_p \rangle\). Therefore
\[
m \cdot\mathbf1_G= m\psi^N - m(\psi^N - 1) \in \mathbf{A}\langle V_p \rangle.
\]
But \(m\) is a rational integer, hence \(m \cdot\mathbf1_G \in \mathbf{R}(G)\). Thus \(m  \cdot\mathbf1_G\in \mathbf{A}\langle V_p \rangle \cap \mathbf{R}(G)=V_p\).
\end{proof}
\begin{thm}[Brauer's Theorem on Induced Characters]\label{Brauer's Theorem}
Every character of irreducible complex representation of \(G\) is a \(\mathbb{Z}\)-linear combination of characters induced from elementary subgroups of \(G\).
\end{thm}
\begin{proof}
Let \(|G| = p^n m\) with \((p,m) = 1\). By Theorem \ref{m in V_p}, we have \(m  \cdot\mathbf1_G\in V_p\). For each \( \chi \in \mathbf{R}(G) \), since \( m \cdot \mathbf{1}_G \in V_p \) and \( V_p \) is an ideal of \( \mathbf{R}(G) \), we have
\[
 m \cdot \chi= m \cdot (\mathbf{1}_G \cdot \chi) = (m \cdot \mathbf{1}_G) \cdot \chi \in V_p.
\]
Thus
\[
m \cdot \mathbf{R}(G) \subseteq V_p.
\]
Then \[
m \cdot \mathbf{R}(G) \le V_p\le\mathbf{R}(G)
\]
as addition abelian groups. Then
\[[\mathbf{R}(G) : V_p]\mid[\mathbf{R}(G) : m\cdot \mathbf{R}(G)]=m^{\operatorname{rank}(\mathbf{R}(G))},\]and then \[p\mid[\mathbf{R}(G) : V_p].\] 

Let \(V = \sum_p V_p\). Then \([\mathbf{R}(G) : V]\mid[\mathbf{R}(G) : V_p],\forall \text{prime }p\mid|G|\). Hence \([\mathbf{R}(G) : V] = 1\), so \(V = \mathbf{R}(G)\). This proves the theorem.
\end{proof}

